\documentclass[12pt,a4paper]{article}

%=========================================================
%                   ENCODAGE ET POLICES
%=========================================================
\usepackage[utf8]{inputenc}
\usepackage[T1]{fontenc}
\usepackage[french]{babel}
\usepackage{lmodern}

%=========================================================
%                     MATHÉMATIQUES
%=========================================================
\usepackage{amsmath,amssymb,amsfonts,mathrsfs}
\usepackage{tkz-tab}
\everymath{\displaystyle}

%=========================================================
%                     MISE EN PAGE
%=========================================================
\usepackage[left=1.5cm,right=1.5cm,top=1.5cm,bottom=1.8cm]{geometry}
\usepackage{multicol}
\usepackage{float}
\usepackage{array,multirow}
\usepackage{enumitem}

%=========================================================
%                GRAPHISMES ET COULEURS
%=========================================================
\usepackage{tikz}
\usetikzlibrary{arrows.meta,calc}
\usepackage{pgfplots}
\pgfplotsset{compat=1.15}

\usepackage[most]{tcolorbox}
\tcbuselibrary{skins,hooks}

%=========================================================
%                  LIENS ET ARRIÈRE-PLAN
%=========================================================
\usepackage[colorlinks=true,linkcolor=blue,urlcolor=blue,citecolor=blue]{hyperref}
\usepackage{eso-pic}

%=========================================================
%                   COULEURS ET STYLES
%=========================================================
\definecolor{myred}{RGB}{180, 40, 40}
\definecolor{myblue}{RGB}{20, 50, 120}
\definecolor{mygreen}{RGB}{30, 120, 40}

%=========================================================
%                       FILIGRANE
%=========================================================
\AddToShipoutPicture{
    \AtTextCenter{
        \makebox[0pt]{
            \rotatebox{45}{
                \textcolor[gray]{0.92}{
                    \fontsize{5cm}{5cm}\selectfont PGB
                }
            }
        }
    }
}

\begin{document}

% En-tête
\begin{center}
    \Large\textbf{\underline{\textcolor{myred}{CHAPITRE 6 : FONCTION LOGARITHME NÉPÉRIEN}}}\\[0.2cm]
    \normalsize\textbf{Prof : M. BA} \hfill \textbf{Classe : Terminale L}\\[0.1cm]
    \textbf{Durée : 6 heures} \hfill \textbf{Année scolaire : 2025 -- 2026}
\end{center}
\hrule
\vspace{0.3cm}

%=========================================================
%   I. Généralités
%=========================================================

\section*{\color{myred}I. Généralités}

\paragraph{\underline{Activité 1}}
Avec la calculatrice remplir le tableau suivant :

\begin{center}
\renewcommand{\arraystretch}{1.5}
\begin{tabular}{|c|c|c|c|c|c|c|c|c|c|}
\hline
$a$ & -2 & -1 & 0 & 0,25 & 0,5 & 0,75 & 1 & 2 & 5 \\
\hline
$\ln a$ & \hspace*{0.6cm} & \hspace*{0.6cm} & \hspace*{0.6cm} & \hspace*{0.6cm} & \hspace*{0.6cm} & \hspace*{0.6cm} & \hspace*{0.6cm} & \hspace*{0.6cm} & \hspace*{0.6cm} \\
\hline
\end{tabular}
\end{center}

\vspace{0.3cm}

\subsection*{\color{myred}1. Définition et notation}
\begin{tcolorbox}[colback=blue!5!white,colframe=myblue,title=\textbf{Définition}]
On appelle \textbf{fonction logarithme népérien}, notée $\ln$, la primitive sur $]0\,;\,+\infty[$ de la fonction $x \mapsto \dfrac{1}{x}$ qui s'annule en $1$.
\[
\ln(1) = 0 \quad \text{et pour tout } x \in ]0\,;\,+\infty[, \quad (\ln x)' = \frac{1}{x}
\]
\end{tcolorbox}

\end{document}